\documentclass[11pt,a4paper]{article}
\usepackage[margin=1in]{geometry}
\usepackage{amsmath,amssymb,amsthm}
\usepackage{booktabs}
\usepackage{xcolor}
\definecolor{linkc}{rgb}{0.0,0.2,0.55}
\usepackage[colorlinks=true,linkcolor=linkc,citecolor=linkc,urlcolor=linkc]{hyperref}
\usepackage{microtype}

\theoremstyle{plain}
\newtheorem{theorem}{Theorem}[section]
\newtheorem{lemma}[theorem]{Lemma}
\newtheorem{proposition}[theorem]{Proposition}
\newtheorem{corollary}[theorem]{Corollary}
\theoremstyle{definition}
\newtheorem{definition}[theorem]{Definition}
\newtheorem{remark}[theorem]{Remark}

\newcommand{\Proved}{\textcolor{linkc}{\textsf{[proved]}}}
\newcommand{\Lean}{\textcolor{linkc}{\textsf{[proved, Lean-verified]}}}
\newcommand{\Meas}{\textcolor{linkc}{\textsf{[measured]}}}
\newcommand{\vv}[1]{v_{#1}}

\title{Rigid phase certificates:\\
deciding nine FRACTRAN holdouts, and a decider that explains why}
\author{%
  \textit{(authorship placeholder --- see \S\ref{sec:data})}%
}
\date{August 2026}

\begin{document}
\maketitle

\begin{abstract}
We decide nine open machines of the FRACTRAN Busy Beaver competition
$\mathrm{BBf}(23)$: each is proved never to halt from its start value
$n=2$. All nine proofs are formalized in Lean~4 without \texttt{mathlib}
and are \texttt{sorry}-free, depending only on Lean's three core axioms.
The proofs are not ad hoc. We isolate a certificate class --- \emph{rigid
phase certificates} --- whose checking is finite, exact and needs no
simulation beyond a bounded entry, prove its soundness, and formalize
that soundness argument in Lean. The nine machines are exactly the
machines our structural detector flags as \emph{rigid} among the $694$
refined $\mathrm{BBf}(23)$ holdouts; the remaining $675$ are
digit-consuming, and we explain why no certificate of this class can
reach them. A second, equivalence decider --- template isomorphism of
certificates --- resolves the nine into six classes, correcting a
plausible-looking but false grouping into three. Along the way the
formalization contradicted our own informal proofs on a point of
principle: checking that a rival rule is disabled at the endpoints of a
run is unsound unless a single guard conjunct is pinned, and we exhibit
a machine-checked counterexample.
\end{abstract}

\section{Introduction}

FRACTRAN \cite{conway1987} is Conway's minimal Turing-complete language.
A program is a finite ordered list of positive rationals; the state is a
positive integer $n$; one step replaces $n$ by $n\cdot p/q$ for the first
listed fraction $p/q$ with $q \mid n$, and the program halts when no
fraction applies. The Busy Beaver variant $\mathrm{BBf}(k)$ asks for the
longest finite run among programs of description size at most $k$
\cite{bbfractran}. Sizes up to $22$ are settled; $\mathrm{BBf}(23)$ is
the live frontier, with $21{,}233$ machines officially undecided as of 10
July 2026 and a refined list of $694$ produced by a decider not yet
independently reproduced.

Deciding a machine on such a list means proving that a specific, fully
explicit program never halts. What makes this hard is that the interesting
survivors read their own digits: the next rule to fire depends on
arithmetic detail that no bounded-state invariant tracks. The exceptions
are machines whose orbit has a \emph{closed form}, and it is a recurring
observation that these fall to analysis while the rest do not.

\paragraph{Contributions.}
\begin{enumerate}
\item \textbf{Nine decisions} (\S\ref{sec:nine}): the FRACTRAN programs at
  lines $13649$, $14793$, $16268$, $18226$, $18380$, $19205$, $20255$,
  $20589$, $20660$ of the official $\mathrm{BBf}(23)$ holdout list never
  halt. \Lean
\item \textbf{A certificate class and a decider} (\S\ref{sec:cert}):
  rigid phase certificates, with a soundness theorem; checking is
  symbolic in $\{1,n,2^n\}$-arithmetic and decides guard positivity for
  \emph{all} phase indices at once. \Proved
\item \textbf{Formalized soundness} (\S\ref{sec:lean}): the soundness
  argument itself is in Lean, and the nine machine theorems are derived
  as instances of one abstract theorem.
\item \textbf{An equivalence decider} (\S\ref{sec:equiv}): template
  isomorphism of certificates, decidable by finite search, which resolves
  the nine into six classes.
\item \textbf{A negative result about endpoint checking}
  (\S\ref{sec:corner}): pinning a single guard conjunct is necessary, not
  merely convenient. \Lean
\item \textbf{A structural census} (\S\ref{sec:scope}) locating the
  boundary of the method: $675$ of the $694$ refined holdouts are outside
  it, by their dynamics rather than by our budget. \Meas
\end{enumerate}

\paragraph{Epistemic conventions.} Every claim carries a label: \Proved{}
for a hand proof, \Lean{} for one machine-checked in Lean~4, \Meas{} for
a measurement over a stated finite domain. We use them because two of the
statements below are measurements that could be mistaken for theorems.

\section{Preliminaries}\label{sec:prelim}

Fix the primes $(2,3,5,7,11)$; every fraction appearing in the machines we
treat factors over them. A state is the exponent vector
$s = (\vv2,\vv3,\vv5,\vv7,\vv{11}) \in \mathbb{N}^5$ of $n$, so that
FRACTRAN's divisibility test becomes a conjunction of coordinate lower
bounds and its update becomes vector addition. A \emph{rule} is a pair
(guard, delta) with guard a conjunction $\bigwedge_i (s_i \ge \theta_i)$;
a \emph{machine} is a priority-ordered list of rules; one step fires the
first rule whose guard holds, and the machine halts when none does. We
write $s \xrightarrow{m} t$ for ``the machine performs $m$ steps from $s$
and is then at $t$'' (in particular it does not halt on the way), and
$\mathrm{NeverHalts}(s)$ for ``every step count is completed from $s$''.

\begin{lemma}[normalized halt criterion]\label{lem:halt}
For each machine, the set of halting states is a conjunction of
coordinate equalities and at most one disjunction, obtained by negating
all five guards and absorbing implications between them. For instance
machine $431$ halts exactly at $\vv2 = \vv7 = 0 \wedge (\vv5 = 0 \vee
\vv{11} = 0)$. \Lean
\end{lemma}

The point of Lemma~\ref{lem:halt} is that non-halting becomes a statement
about which coordinates stay positive, which is what certificates track.

\section{Rigid phase certificates}\label{sec:cert}

\begin{definition}[certificate]\label{def:cert}
Let $\mathrm{EXP} = \{\,a + bn + c\,2^{n} : a,b,c \in \mathbb{Q}\,\}$. A
\emph{rigid phase certificate} for a machine $M$ consists of
\begin{itemize}
\item an \emph{entry} $(E,i_0)$: from the start state, $M$ reaches
  $B(i_0)$ in $E$ steps;
\item a \emph{boundary family} $B(n)$, each coordinate in $\mathrm{EXP}$;
\item \emph{branches}, one per residue class of $n$ modulo $2$ (or a
  single branch), each a list of \emph{stages}:
  a \emph{run} $(\text{rule}, \text{count})$ with count in $\mathrm{EXP}$,
  or a \emph{block} $(K, w, \sigma)$ with $K \in \mathrm{EXP}$, $w$ a
  fixed word of runs, and $\sigma(k)$ a block-start state affine in the
  block index $k$ with $\mathrm{EXP}$ constant terms and integer slopes.
\end{itemize}
\end{definition}

The checker verifies five conditions. \textbf{C1 (chain)}: the symbolic
composition of the stages equals $B(n+1)$, and each block's word-delta
equals its start function's $k$-slope. Because $\{1,n,2^n\}$ are linearly
independent over $\mathbb{Q}$, state equality is a finite comparison of
coefficients. \textbf{C2 (guards)}: the acting rule's guard holds at the
corners of each block rectangle. \textbf{C3 (priority)}: for each
higher-priority rule, one \emph{fixed} guard conjunct fails at all
corners. \textbf{C4 (entry)}: by direct simulation. \textbf{C5
(coverage)}: the branches cover every $n \ge i_0$.

Two facts make this a decision procedure rather than a search.

\begin{proposition}[decidable eventual positivity]\label{prop:evpos}
For $g \in \mathrm{EXP}$, a threshold $t$, and an arithmetic progression
of indices, ``$g(n) \ge t$ for all admissible $n \ge n_0$'' is decidable.
\Proved
\end{proposition}
\begin{proof}
If the $2^n$-coefficient is negative, or is zero and the linear
coefficient is negative, $g$ is eventually below any threshold and the
answer is no. Otherwise $g$ is eventually nondecreasing along the
progression, with an explicitly computable index $N$ past which its
increment is nonnegative; check $n_0 \le n \le N$ directly.
\end{proof}

\begin{theorem}[soundness]\label{thm:sound}
If the checker accepts a certificate for $M$, then $M$ never halts from
its start state. \Lean
\end{theorem}

The proof factors into a purely dynamical part, which is what we
formalize, and the bookkeeping that C1--C5 discharge.

\begin{theorem}[phase induction; the certificate contract]\label{thm:A}
Let $B : \mathbb{N} \to \mathbb{N}^5$, and suppose
$s_0 \xrightarrow{E} B(i_0)$, and that for every $i \ge i_0$ there is
$m \ge 1$ with $B(i) \xrightarrow{m} B(i+1)$. Then
$\mathrm{NeverHalts}(s_0)$. \Lean
\end{theorem}
\begin{proof}
By induction, for every $k$ there are $m \ge k$ and a state $t$ with
$s_0 \xrightarrow{m} t$: the base case is the entry, and each phase adds
at least one step. A machine that completes arbitrarily many steps
completes every step count.
\end{proof}

C4 supplies Theorem~\ref{thm:A}'s entry hypothesis; C1--C3 supply its
phase hypothesis for the branch containing $i$; C5 ensures every
$i \ge i_0$ has one. That the phase hypothesis really follows from
C1--C3 is where the geometry enters, and it is the subject of the next
section.

\section{Corners, and why one conjunct must be pinned}\label{sec:corner}

Within a run the state is affine in the run counter; across a block the
start state is affine in the block index. Hence every guard value inside
a block has the form $a + bk + ct$ with $(k,t)$ ranging over a rectangle.

\begin{theorem}[corner principle]\label{thm:B}
If $a + bk + ct \ge 0$ at the four corners of $[0,K]\times[0,L]$, then
$a+bk+ct \ge 0$ on the whole rectangle; likewise for $\le \theta$.
\Lean
\end{theorem}
\begin{proof}
A linear integer function attains its extremes on an interval at an
endpoint; apply this in $k$ along the two extreme rows, then in $t$ for
each fixed $k$.
\end{proof}

Theorem~\ref{thm:B} licenses C2 immediately: finitely many corner tests
replace infinitely many pointwise tests. C3 is subtler, and here our
first informal treatment was wrong in a way worth recording.

A rule is \emph{disabled} when \emph{some} conjunct of its guard fails.
It is tempting to verify ``disabled'' at the endpoints of a run and
conclude it holds throughout. That inference is invalid: the endpoints
may be witnessed by \emph{different} conjuncts.

\begin{proposition}[endpoint testing of a disjunction is unsound]\label{prop:unsound}
There are affine conjunct-values $p+ut$ and $q+vt$ and a run length $n$
such that at $t=0$ and at $t=n$ some conjunct fails, yet at an
intermediate $t$ both hold --- so the rule fires in the middle of a run
the check has accepted. \Lean
\end{proposition}
\begin{proof}
Take $p+ut = t$ and $q+vt = 2-t$ with $n = 2$. At $t=0$ the first value
is $0$; at $t=2$ the second is $0$; at $t=1$ both equal $1$.
\end{proof}

Accordingly C3 requires a \emph{single fixed} conjunct to fail at all
corners, which is sound by Theorem~\ref{thm:B}. The same discipline
appears in the Lean development as a design constraint: each firing lemma
names the coordinate that disables each higher-priority rival, and a rule
typically needs \emph{two} such lemmas because different coordinates do
the work in different stages --- machine $680$'s rule $25/2$ is protected
by $\vv5 = 0$ inside its blocks but by $\vv3 = \vv7 = 0$ during its
opening run. This was the single most common source of error in the
formalization, and it is exactly the content of C3.

\section{The nine machines}\label{sec:nine}

\begin{theorem}\label{thm:nine}
Each program in Table~\ref{tab:nine}, started at $n=2$, never halts.
\Lean
\end{theorem}

\begin{table}[t]
\centering\small
\begin{tabular}{llll}
\toprule
line & program & boundary family & phase length\\
\midrule
$13649$ & $[5/6,\,9/35,\,8/55,\,7/2,\,605/7]$ & $(2^{i+1}\!-\!1,0,0,0,i)$ & $2^{i+3}-5$\\
$14793$ & $[63/10,\,8/77,\,33/2,\,5/9,\,7/3]$ & $(2^i\!+\!1,2^i\!-\!2,0,0,i\!-\!1)$ & $4\cdot 2^{i}$\\
$16268$ & $[9/70,\,25/2,\,44/15,\,7/55,\,3/5]$ & $(0,0,2^{i+2},0,2^{i+1}\!-\!1)$ & $7(2^{i+1}\!-\!1)+4$\\
$18226$ & $[9/35,\,5/6,\,8/55,\,7/2,\,605/7]$ & transport of line $13649$ & $2^{i+3}-5$\\
$18380$ & $[7/15,\,9/14,\,125/77,\,2/5,\,847/2]$ & transport of line $13649$ & $2^{i+3}-5$\\
$19205$ & $[9/10,\,5/21,\,343/55,\,2/7,\,605/2]$ & transport of line $13649$ & $2^{i+3}-5$\\
$20255$ & $[8/15,\,147/22,\,35/2,\,11/49,\,3/7]$ & $(0,0,2^k\!+\!k,2^{k+1}\!-\!1,0)$ & $5\cdot 2^{k}$\\
$20589$ & $[77/30,\,88/21,\,9/2,\,5/11,\,7/3]$ & $(0,2^{m+1}\!+\!2m,0,0,2^m\!-\!1)$ & $5\cdot 2^{m}$\\
$20660$ & $[9/70,\,44/15,\,25/2,\,7/55,\,3/5]$ & $(2^m,0,0,0,2^m\!-\!1)$ & $6\cdot 2^{m}-3$\\
\bottomrule
\end{tabular}
\caption{The nine decided machines, by line number in the official
$21{,}233$-machine $\mathrm{BBf}(23)$ holdout list of 10 July 2026. All
nine also appear in the refined $694$-machine list.}
\label{tab:nine}
\end{table}

\paragraph{A worked phase.} For line $13649$, write $W = 2^{i+1}-1$ and
$m = 2^i-1$. From $B_i = (W,0,0,0,i)$ the guards force
\[
  f_3^{\,W},\quad (f_4 f_1)^{m} f_4,\quad f_2,\quad
  \begin{cases}
  (f_0^3 f_2)^{q}\, f_2^{\,2q}, & i \text{ even},\ 2m = 3q,\\
  (f_0^3 f_2)^{q}\, f_0^2\, f_2^{\,2q+2}, & i \text{ odd},\ 2m = 3q+2,
  \end{cases}
\]
landing exactly on $B_{i+1}$ in $2^{i+3}-5$ steps. The case split is
forced by $2^{i+1} \bmod 3$; the only nontrivial guard in the entire
phase is $\vv{11} \ge 1$ at each $f_2$, and $\vv{11}$ has floor exactly
$i$ over the phase. That floor is the machine's safety margin, and it
grows linearly.

\paragraph{Three mechanisms, not one.} Lines $13649$ and $14793$ survive
by a \emph{linearly growing fuel} ($\vv{11} = i$, resp.\ $i-1$). Line
$16268$ survives by an \emph{exactly preserved invariant}: its boundary
satisfies $\vv5 = 2(\vv{11}+1)$ and the phase map sends
$(\vv5,\vv{11}) \mapsto (2\vv5, 2\vv{11}+1)$, which preserves it. Line
$20255$ is the only machine of the nine with \emph{no} parity split ---
its blocks consume nothing three at a time, so no residue condition
arises. Lines $20589$ and $20660$ have \emph{rotated} phase words: the
same multiset of rules per block in a different cyclic order, with the
rotation depending on the parity.

\paragraph{Margins are not thin.} It is worth contrasting these with the
one previously decided machine of comparable shape, a $\mathrm{BB}(6)$
holdout settled in April 2026 whose halting condition reduces to a
$2$-adic valuation of $2\cdot 3^i + i + 5$ staying above a linear bound.
That margin is \emph{thin}: it can only be controlled by transcendence
methods, and the published proof invokes Baker--W\"ustholz. None of our
nine has a thin margin, and no result beyond guard arithmetic and
induction is used anywhere below.

\section{The equivalence decider}\label{sec:equiv}

Certificates support a second decision problem. Say two certified
machines are \emph{template-isomorphic} if some bijection of rule indices
and permutation of the prime axes carries one certificate exactly onto
the other. This is decidable by finite search over at most $5!\cdot 5!$
maps, and an isomorphism yields an exact correspondence of orbits from
the entry boundaries onward.

\begin{theorem}\label{thm:equiv}
Under template isomorphism the nine machines fall into six classes:
$\{13649, 18226, 18380, 19205\}$ and five singletons. \Proved
\end{theorem}

This corrects a natural error. Inspection suggests three families, since
$20255$, $20589$, $20660$ each share a fraction multiset with an earlier
machine up to reordering. But reordering fractions changes rule
\emph{priority}, and priority changes the firing word: those three have
genuinely different phase templates, which is why each needed its own
proof. The decider settles the question formally rather than by
inspection --- and the Lean development independently reproduces the same
boundary (\S\ref{sec:lean}).

\section{The Lean development}\label{sec:lean}

The formalization is in Lean~4 with \textbf{no} \texttt{mathlib}: only
core arithmetic, \texttt{omega} and kernel \texttt{decide}. It comprises
$3{,}262$ lines in $9$ files, builds from scratch in about four seconds,
contains no \texttt{sorry}, no \texttt{admit}, no \texttt{native\_decide}
and no \texttt{axiom} declarations, and each of the nine non-halting
theorems depends on exactly Lean's three core axioms
(\texttt{propext}, \texttt{Classical.choice}, \texttt{Quot.sound}); the
entry lemmas depend on none.

\texttt{Fractran.lean} defines the semantics and proves the reusable
core: a block-iteration lemma and the affine endpoint lemmas.
\texttt{Decider.lean} proves Theorem~\ref{thm:A}, Theorem~\ref{thm:B} and
Proposition~\ref{prop:unsound}, and re-derives four of the nine machine
theorems through Theorem~\ref{thm:A}, confirming that the certificate
contract is exactly what the machine proofs establish.

Two byproducts deserve mention.

\paragraph{Transport as a proof principle.} Three of the nine are
transports of line $13649$. Rather than copy its development, we prove a
\emph{firing specification}: six transported firing obligations, stated on
the source machine's role coordinates and embedded by an axis map, from
which the entire phase development and induction follow generically.
Instantiating it requires discharging only those six facts against the
target machine's own rule list and priority order. The specification
demands the top-priority rule fire only at states with $\vv7 = 0$, which
is precisely what makes the priority swaps in two of the transports
harmless --- and precisely what the rotated machines cannot satisfy. The
formal development thus rediscovers the equivalence boundary of
Theorem~\ref{thm:equiv} from the other side.

\paragraph{Entries collapse.} Our informal case files entered the
induction at boundary index $4$ after $36$, $76$ or $78$ steps, because
their phase templates had been stated only for large indices. The Lean
phase lemmas hold for \emph{all} parameter values, and the true entries
are far shorter: one step for lines $20255$ and $20589$, sixteen for
$16268$, and \emph{zero} for line $20660$, whose start state $(1,0,0,0,0)$
already \emph{is} its index-$0$ boundary. The old counts survive as
ground-truth cross-checks, which is how we confirmed the re-indexing was
faithful rather than a slip.

\section{Scope: what rigid certificates cannot do}\label{sec:scope}

A detector for the certificate class costs little: accelerate the machine
exactly, segment its run at occurrences of a marker rule, and classify the
resulting per-phase series. We ran it over both holdout lists.

\begin{table}[h]
\centering\small
\begin{tabular}{lrrr}
\toprule
list & rigid (GEO) & polynomial & non-rigid\\
\midrule
refined, $694$ machines & $9$ & $10$ & $675$\\
official, $21{,}233$ machines & $3{,}253$ & $6{,}797$ & $11{,}183$\\
\bottomrule
\end{tabular}
\caption{Structural census of the $\mathrm{BBf}(23)$ holdouts. \Meas}
\label{tab:census}
\end{table}

The reading is that rigidity is the \emph{decidable margin} of the
holdout set: about $15\%$ of the official list is rigid, and after the
community's newest deciders have run, exactly $9$ rigid machines survive
--- the nine decided here. What remains is overwhelmingly non-rigid:
$675$ of $694$. Those machines choose their next rule from arithmetic
detail that grows with the state, so no boundary family in $\mathrm{EXP}$
--- indeed no family with a bounded description --- tracks them, and the
certificate class is empty for them by construction rather than by
budget. This is the same phenomenon that makes generalized Collatz
problems hard, and we make no claim to address it.

We also note what did \emph{not} appear: among the nine rigid survivors,
none has a thin, Baker-type margin. If machines of that kind exist on
these lists, they are not in the rigid class we can detect, and the
natural place to look is the $\mathrm{BB}(6)$ Turing-machine list, where
extracting an arithmetic reduction is itself the hard step.

\section{Data and reproducibility}\label{sec:data}

All code, certificates, logs and the Lean development accompany this
paper. The checker runs the nine certificates in about a second; a
falsifier that systematically perturbs boundaries, counts, rule indices
and block coefficients produced $540$ corrupted certificates, all
rejected, none accepted. Every machine proof is additionally cross-checked
against direct big-integer FRACTRAN simulation.

\paragraph{One faithfulness boundary, stated plainly.} The Lean
development formalizes the exponent-vector machine, and its agreement
with divisibility FRACTRAN --- that dividing by $q$ is subtracting $q$'s
exponent vector --- is checked computationally rather than proved, since
a formal proof needs unique factorization. Closing that gap is
straightforward with \texttt{mathlib} and is the obvious next step for
anyone wishing to remove the last informal link.

\paragraph{Provenance.} Holdout lists are those published by the
$\mathrm{BBf}$ project: the official size-$23$ list of 10 July 2026
($21{,}233$ machines) and the refined list of 25 July 2026 ($694$
machines, produced by a decider not yet independently reproduced; we do
not rely on it --- all nine machines are on the official list).

\begin{thebibliography}{9}
\bibitem{conway1987} J.~H. Conway. Fractran: a simple universal
programming language for arithmetic. In \textit{Open Problems in
Communication and Computation}, Springer, 1987, 4--26.
\bibitem{bbfractran} The BBFractran project. Holdout lists and
enumeration conventions for $\mathrm{BBf}(k)$.
\url{https://github.com/int-y1/BBFractran}
\bibitem{bbchallenge} The bbchallenge Collaboration.
\url{https://wiki.bbchallenge.org}
\bibitem{lean4} L. de Moura and S. Ullrich. The Lean 4 theorem prover and
programming language. \textit{CADE-28}, 2021.
\end{thebibliography}

\end{document}
